\section{Plano de despliegue: pila tecnológica y coordinación organizacional}

\subsection{Pila tecnológica de referencia}
Combinando los puntos de vista del curso, una pila de agentes viable puede dividirse en seis capas: capa de interfaz, capa de orquestación, capa de memoria, capa de herramientas, capa de validación y capa de observabilidad. Cada capa debe contar con una capacidad mínima de autonomía y comunicarse mediante un protocolo estable.

\begin{table}[H]
\centering
\begin{tabular}{p{2.4cm}p{4.0cm}p{4.6cm}}
\toprule
\textbf{Capa} & \textbf{Responsabilidad principal} & \textbf{Producto clave} \\
\midrule
Capa de interfaz & Recibir tareas, verificar permisos, dar formato a las salidas & Task spec, contexto de permisos \\
Capa de orquestación & Descomponer tareas, enrutar, disparar la replanificación & Action graph, plan de ejecución \\
Capa de memoria & Lectura y escritura por niveles, compresión, descarte, resolución de conflictos & Índice y resumen de memoria \\
Capa de herramientas & Protocolo unificado de herramientas, encapsulado de errores, reintentos & Catálogo de API invocables \\
Capa de validación & Validación de pasos y de resultado final, intercepción de riesgos & Reporte de validación y etiquetas de riesgo \\
Capa de observabilidad & Rastreo de trayectorias, agregación de métricas, alertas & Trace, dashboard, bitácora de auditoría \\
\bottomrule
\end{tabular}
\caption{Estructura de referencia de seis capas para el despliegue de sistemas de agentes}
\end{table}

\subsection{Modelo de coordinación organizacional}
Un problema organizacional común en los proyectos de agentes es que el grupo de modelos, el grupo de plataforma y el grupo de negocio optimizan cada uno sus propios objetivos parciales. La lección del curso es establecer la «tasa de éxito de la tarea» como métrica guía unificada y hacer que los distintos equipos compartan el mismo conjunto de datos de trayectorias y la misma clasificación de fallas.

\begin{importantbox}{Tres prácticas para la coordinación organizacional}
\begin{enumerate}
    \item Unificar el esquema de eventos para evitar criterios distintos entre grupos;
    \item Unificar el mecanismo de reversión para mantener controlado el riesgo en producción;
    \item Unificar el marco de prioridades, corrigiendo primero las fallas de mayor pérdida.
\end{enumerate}
\end{importantbox}

\subsection{Gobernanza y cumplimiento}
En los despliegues empresariales, el agente no es solo un problema técnico: también involucra la minimización de permisos, el manejo de datos sensibles, el registro de auditoría y la atribución de responsabilidades. El curso recomienda integrar la capacidad de gobernanza dentro del sistema desde el diseño, en lugar de aplicarla como un parche de cumplimiento después del lanzamiento.

\begin{warningbox}{La capacidad de gobernanza debe establecerse por adelantado}
\begin{itemize}
    \item Sin una jerarquía de permisos, la invocación de herramientas puede generar riesgos de exceso de privilegios;
    \item Sin una bitácora de auditoría, es difícil rastrear la responsabilidad después de un incidente;
    \item Sin reglas de límite estricto, una anomalía del modelo puede propagarse rápidamente.
\end{itemize}
\end{warningbox}

\begin{knowledgebox}{Lista de verificación previa al despliegue}
\begin{itemize}
    \item ¿Se cuenta con reversión a nivel de tarea y con la posibilidad de intervención humana?
    \item ¿Se cuenta con registros a nivel de paso y con trayectorias que puedan reproducirse?
    \item ¿Se cuenta con límites de presupuesto y con mecanismos de corte ante fallas?
    \item ¿Se completaron pruebas de equipo rojo para los escenarios sensibles?
\end{itemize}
\end{knowledgebox}

\subsection{Resumen de la sección}
El despliegue de agentes es una ingeniería conjunta de «tecnología + organización + gobernanza». Solo diseñando los tres elementos de manera integrada puede el sistema atravesar de forma estable la brecha entre la demostración y la producción.
